\subsubsection{Mixed surface boundaries with formally linear exceptional curves}

\begin{theorem}[Proved under stated hypotheses: mixed finite boundary with formal cones]
\label{D:formal-cone-mixed-boundary}
Let $k$ be algebraically closed of characteristic $p>0$, let
$S$ be a normal connected proper surface over $k$, and let
$Y\subset S$ be finite.  Suppose that every singular point
$s\in Y$ satisfies at least one of the following conditions:
\begin{enumerate}
\item $S$ is $\mathbb Q$-factorial at $s$;
\item $S$ is log canonical at $s$;
\item a resolution of the germ at $s$ has reduced exceptional
      fibre one smooth curve $C$, and its formal neighborhood
      is isomorphic to the formal zero section of its negative
      normal line bundle $N=N_{C/\widetilde S}$.
\end{enumerate}
Then $\pi_1(S\setminus Y)$ is topologically finitely generated.
The singularities retained in the open are unrestricted.

In condition (3), it is enough that the normal bundle satisfies
\[
 H^1(C,T_C\otimes N^{-n})=0\quad(n\geq1),
 \qquad
 H^1(C,N^{1-n})=0\quad(n\geq2).
\]
In particular, if $C^2=-d<0$ and $g=g(C)$, it is enough that
$d>\max\{0,4g-4\}$.  This includes every single smooth
elliptic exceptional curve.

For a possibly nonnormal proper scheme of dimension at most
two, the same finite-generation conclusion holds if $Y$ has
codimension at least two in every irreducible component and
its finite inverse image on each normal surface component
satisfies the displayed alternatives.
\end{theorem}
\begin{proof}
The cohomological and numerical assertions follow from
Lemma~\ref{D:negative-curve-linearization}.  Thus it suffices
to prove the theorem under alternatives (1)--(3).

Let $T\subset Y$ consist of the singular points which are
$\mathbb Q$-factorial.  We leave these points unchanged.  At
each remaining point of $Y$, choose the local resolution in
(3), or the minimal resolution in (2).  In the latter case,
Lemma~\ref{D:lc-exceptional-configurations} gives either a
rational tree, one smooth elliptic curve, or a rational nodal
cycle.  The rational-tree case would be $\mathbb Q$-factorial
by Proposition~\ref{D:rational-tree-Qfactorial}, so it already
belongs to $T$.  Blow up a node when necessary to give the
cycle smooth components.  At singular points of $S\setminus Y$
choose any regular resolution, and keep regular points
unchanged.

These local choices produce a proper birational map
$q:X\to S$.  Explicitly, choose a neighborhood of each of the
finitely many singular points containing no other singular
point and spread the specified local model to a smaller such
neighborhood.  Each chosen resolution is an isomorphism on
the punctured neighborhood, so it glues with the identity on
the regular locus and with the other resolutions.  Properness
is local on the target.  The resulting $X$ is regular away
from the retained $\mathbb Q$-factorial points and is therefore
$\mathbb Q$-factorial everywhere.  It is proper, hence
projective by the intersection argument in
Lemma~\ref{D:Qfactorial-partial-resolution}.  That argument
derives projectivity from properness and $\mathbb Q$-factoriality;
it does not require the prescribed local resolution to have
been chosen projective in advance.

Put $U=S\setminus Y$ and $U'=q^{-1}(U)$.  If $V\to U$ is a
connected finite \'etale cover, its pullback to the normal
integral $U'$ is connected: every finite \'etale component
dominates $U'$, whereas its generic algebra is the single
field $k(V)$.  Thus $\pi_1(U')\to\pi_1(U)$ is surjective.
The finite normalization of $S$ in $k(V)$ is \'etale away from
$Y$, so at any omitted point its restriction to the punctured
germ is \'etale.

Every divisorial component of $X\setminus U'$ lies over one
of the points of $Y\setminus T$ which we resolved.  At a
formally linear curve, Proposition~\ref{D:graded-cone-tame}
proves tameness at its valuation, including separability of
the residue extension.  At an lc elliptic or cyclic core,
Corollary~\ref{D:lc-core-valuative-tame} proves tameness, also at
the node subdivisions.  Retained points of $T$, and regular
omitted points, have codimension two in $X$ and create no
divisorial boundary condition.  Exceptional curves over
singularities retained in $U$ are contained in $U'$.
Consequently every finite quotient of $\pi_1(U)$ factors
through the fixed-boundary tame quotient of $\pi_1(U')$, and
we have a surjection
\[
                  Q(X,X\setminus U')\longrightarrow\pi_1(U).
\]
Its source is topologically finitely generated by
Proposition~\ref{D:fixed-boundary-positive}, because $X$ is
projective.  This proves the normal case.

For the final assertion, pass to the finite disjoint union
of the normalizations of the reduced irreducible components.
The proved result applies to every normal surface component.
Components of dimension at most one have no omitted points
by the codimension hypothesis, so they are proper and their
fundamental groups are finitely generated.  The finite proper
descent argument in Proposition~\ref{D:fixed-boundary-positive},
using \cite[IX, Corollary 5.2]{SGA1}, gives finite generation
downstairs.
\end{proof}

\begin{corollary}[Proved under stated hypotheses: fixed boundary with isolated formal cones]
\label{D:formal-cone-fixed-boundary}
Let $S$ be a normal connected proper surface over an algebraically
closed field of characteristic $p>0$, and let $Y\subsetneq S$
be any closed subset.  Assume that every singular point
$s\in Y$ which is not $\mathbb Q$-factorial is isolated in
$Y$ and is either log canonical or satisfies condition (3)
of Theorem~\ref{D:formal-cone-mixed-boundary}.  Then
\[
                         Q(S,Y)
\]
is topologically finitely generated.  The cohomological and
numerical sufficient conditions in that theorem apply here
as well.  Singularities outside $Y$ are unrestricted.
\end{corollary}
\begin{proof}
As in the preceding proof, retain every $\mathbb Q$-factorial
singular point on $Y$, resolve the other singular points with
the indicated models, and resolve all singular points outside
$Y$ arbitrarily.  The result is a normal proper
$\mathbb Q$-factorial, hence projective, surface $q:X\to S$.
Put $U=S\setminus Y$ and $U'=q^{-1}(U)$.  The birational
pullback again gives a surjection $\pi_1(U')\to\pi_1(U)$.

Consider a connected finite \'etale cover of $U$ which is
tame at the original codimension-one boundary valuations.
Their strict transforms on $X$ define the same DVRs, so the
pulled-back cover remains tame there.  Every new exceptional
boundary curve lies over a non-$\mathbb Q$-factorial singular
point $s\in Y$.  By hypothesis $s$ is isolated in $Y$; on a
sufficiently small neighborhood the cover is therefore
\'etale on the punctured germ.  Its ramification at the new
exceptional valuations is tame by
Proposition~\ref{D:graded-cone-tame} or
Corollary~\ref{D:lc-core-valuative-tame}, exactly as in the preceding
proof.  Resolutions above points outside $Y$ remain inside
$U'$.  Thus pullback defines a surjection
\[
                         Q(X,X\setminus U')\longrightarrow Q(S,Y).
\]
The projective fixed-boundary theorem makes the source
topologically finitely generated.  This proves the assertion
without dropping the isolation condition on the newly
resolved boundary points.
\end{proof}
